\documentclass[12pt,a4paper,oneside]{report}

\usepackage[utf8]{inputenc}
\usepackage[T1]{fontenc}
\usepackage[italian]{babel}
\usepackage[a4paper,margin=2.5cm]{geometry}
\usepackage{graphicx}
\usepackage{booktabs}
\usepackage{siunitx}
\usepackage[hidelinks]{hyperref}

\title{Analisi FEM di un piede bionico}
\author{Amos Rafael Grusovin}
\author{Marco Libertone}
\date{Anno accademico 2026--2027}

\begin{document}

\begin{titlepage}
    \centering
    {\Large Università degli studi di Modena e Reggio-Emilia\par}
    \vspace{1.5cm}
    {\large Bioingegneria per l'innovazione in medicina\par}
    \vspace{2.5cm}
    {\Huge\bfseries Analisi FEM di un piede bionico\par}
    \vspace{0.5cm}
    {\Large Eventuale sottotitolo\par}
    \vfill
    \begin{flushleft}
        \textbf{Docente:} Edoardo Artioli\\[0.3cm]
        \textbf{Studente:} Amos Rafael Grusovin\\
        \textbf{Studente:} Marco Libertone
    \end{flushleft}
    \vfill
    {\large Anno accademico 2026--2027\par}
\end{titlepage}

\pagenumbering{roman}

\chapter*{Abstract}
\addcontentsline{toc}{chapter}{Abstract}
Riassumere qui l'argomento, gli obiettivi, il metodo seguito e i risultati
principali della tesina. L'abstract dovrebbe essere breve e comprensibile
anche senza aver letto il resto del documento.

Il progetto consiste nel realizzare un'analisi FEM (Finenite Element Method) di un piedi bionico.
Lo studio si basa su una tesi di laurea magistrale di J. Tabucol et al. dal quale 
abbiamo estrapolato il modello 3D della protesi \textit{MyFlex-$\gamma$ foot} in maniera occhiometrica-proporzionale
rispetto alle misure indicate nell'articolo. Una volta ottenuto il modello 2D e 3D 
abbiamo proceduto alla meshatura e alle successive analisi FEM.


\tableofcontents
\clearpage
\renewcommand{\thepage}{\arabic{page}}

\chapter{Introduzione}
\section{Contesto}
Presentare il tema e il contesto in cui si inserisce il lavoro.

\section{Obiettivi e struttura della tesina}
Descrivere gli obiettivi e anticipare brevemente l'organizzazione dei capitoli.

\chapter{Stato dell'arte}
Riassumere le conoscenze e le soluzioni esistenti, citando le fonti utilizzate
come nell'esempio \cite{riferimento-esempio}.

\chapter{Progettazione e metodologia}
\section{Requisiti}
Elencare i requisiti funzionali e i vincoli del progetto.

\section{Metodo di lavoro}
Descrivere gli strumenti, le ipotesi e le fasi seguite.

\chapter{Modellazione e analisi}
\section{Modello}
Descrivere il modello sviluppato e le sue principali caratteristiche.

\section{Verifiche}
Riportare le verifiche svolte, le condizioni al contorno e le ipotesi adottate.
Per esempio, una grandezza può essere espressa con unità di misura:
\[
    F = m\,a \qquad [F] = \si{\newton}.
\]

\chapter{Risultati}
Presentare i risultati in modo ordinato e commentarli.

\begin{table}[htbp]
    \centering
    \caption{Esempio di tabella dei risultati.}
    \label{tab:risultati}
    \begin{tabular}{lcc}
        \toprule
        Caso & Valore iniziale & Valore finale \\
        \midrule
        Caso A & -- & -- \\
        Caso B & -- & -- \\
        \bottomrule
    \end{tabular}
\end{table}

\chapter{Discussione}
Interpretare i risultati, confrontarli con gli obiettivi e indicare eventuali
limiti del lavoro.

\chapter{Conclusioni e sviluppi futuri}
Riassumere i risultati principali e proporre possibili sviluppi futuri.

\appendix
\chapter{Materiale supplementare}
Inserire, se necessario, dettagli tecnici, calcoli o dati aggiuntivi.

\begin{thebibliography}{9}
    \bibitem{riferimento-esempio}
    Cognome, N., \emph{Titolo del testo o dell'articolo},
    Editore o rivista, anno.
\end{thebibliography}

\end{document}